% lp2graph canonical LaTeX
% Reversible with lp2graph.codec.from_canonical_latex (schema 0.1.0).
%@ meta id=hsr_trp family=milp schema=0.1.0
%@ name :: Timetable Rescheduling after Earthquake (highspeed railway)
%@ desc :: MILP model for post‑earthquake timetable rescheduling that captures track outages, track and station recovery, and dynamic passenger demand. Binary variables indicate whether a train service is scheduled, whether a track segment or station is operational, and continuous variables measure unserved demand and delay. The objective minimizes a weighted sum of unmet demand and delay while respecting infrastructure availability.
%@ index P ordered=0 cyclic=0 :: Origin‑destination (OD) pairs.
%@ index T ordered=0 cyclic=0 :: Discrete time periods in the recovery horizon.
%@ index L ordered=0 cyclic=0 :: Track segments.
%@ index S ordered=0 cyclic=0 :: Stations.
%@ param demand shape=P,T kind=matrix domain=passenger_volume :: Dynamic passenger demand for OD pair p in period t after the disruption.
%@ param diff shape=P,T kind=matrix domain=time :: Time difference = time_t - original_time_p.
%@ param seg shape=P,L kind=matrix domain=bool :: Indicator =1 if track segment l belongs to the path of OD pair p, otherwise 0.
%@ param sta shape=P,S kind=matrix domain=bool :: Indicator =1 if station s lies on the path of OD pair p, otherwise 0.
%@ param init_track shape=L kind=vector domain=bool :: Initial operational status of each track segment (0 = out, 1 = functional) at the beginning of the horizon.
%@ param init_station shape=S kind=vector domain=bool :: Initial operational status of each station at the beginning of the horizon.
%@ param max_demand shape=- kind=scalar domain=passenger_volume :: Upper bound on passenger demand used for big‑M linearization.
%@ param M_delay shape=- kind=scalar domain=time :: Big‑M constant for delay linearization.
%@ param alpha shape=- kind=scalar domain=cost :: Weight for unserved demand in the objective.
%@ param beta shape=- kind=scalar domain=cost :: Weight for delay in the objective.
%@ param prev shape=T,T kind=matrix domain=bool :: 1 if the first index is the immediate predecessor of the second index, otherwise 0.
%@ var x shape=P,T domain=binary role=primary drole=- lo=- hi=- :: 1 if a train serving OD pair p is scheduled in period t.
%@ var y shape=L,T domain=binary role=primary drole=- lo=- hi=- :: 1 if track segment l is operational (recovered) in period t.
%@ var z shape=S,T domain=binary role=primary drole=- lo=- hi=- :: 1 if station s is operational (recovered) in period t.
%@ var u shape=P,T domain=non_negative role=primary drole=- lo=0 hi=- :: Unserved passenger demand for OD pair p in period t.
%@ var delta shape=P,T domain=non_negative role=primary drole=- lo=0 hi=- :: Delay (in time units) of the scheduled service for OD pair p in period t relative to the original timetable.
%@ obj sense=min name=resilience_score combination=sum :: Minimize weighted sum of unmet demand and service delay.
%@ con service_track kind=linear domain=- indicator=- :: Service can be scheduled only if every required track segment is operational.
%@ con service_station kind=linear domain=- indicator=- :: Service can be scheduled only if every required station is operational.
%@ con unserved_demand kind=linear domain=- indicator=- :: Definition of unserved demand using big‑M linearization.
%@ con service_delay kind=linear domain=- indicator=- :: Definition of service delay using big‑M linearization.
%@ con track_monotonic kind=linear domain=- indicator=- :: Once a track segment is recovered it stays recovered.
%@ con station_monotonic kind=linear domain=- indicator=- :: Once a station is recovered it stays recovered.
%@ con init_track_status kind=linear domain=- indicator=- :: Initial status of each track segment.
%@ con init_station_status kind=linear domain=- indicator=- :: Initial status of each station.
\begin{align}
  \min\quad & \sum_{p \in \mathcal{P}, t \in \mathcal{T}} alpha \cdot u_{p,t} + \sum_{p \in \mathcal{P}, t \in \mathcal{T}} beta \cdot delta_{p,t} \tag{resilience_score} \\
  & x_{p,t} \le y_{l,t} + 1 - seg_{p,l} \qquad \forall p \in \mathcal{P}, t \in \mathcal{T}, l \in \mathcal{L} \tag{service_track} \\
  & x_{p,t} \le z_{s,t} + 1 - sta_{p,s} \qquad \forall p \in \mathcal{P}, t \in \mathcal{T}, s \in \mathcal{S} \tag{service_station} \\
  & u_{p,t} \ge demand_{p,t} - max_demand + max_demand \cdot x_{p,t} \qquad \forall p \in \mathcal{P}, t \in \mathcal{T} \tag{unserved_demand} \\
  & delta_{p,t} \ge diff_{p,t} - M_delay + M_delay \cdot x_{p,t} \qquad \forall p \in \mathcal{P}, t \in \mathcal{T} \tag{service_delay} \\
  & y_{l,t} \ge \sum_{tp \in \mathcal{T}} prev_{tp,t} \cdot y_{l,tp} \qquad \forall l \in \mathcal{L}, t \in \mathcal{T} \tag{track_monotonic} \\
  & z_{s,t} \ge \sum_{tp \in \mathcal{T}} prev_{tp,t} \cdot z_{s,tp} \qquad \forall s \in \mathcal{S}, t \in \mathcal{T} \tag{station_monotonic} \\
  & y_{l,t} \ge init_track_{l} \qquad \forall l \in \mathcal{L}, t \in \mathcal{T} \tag{init_track_status} \\
  & z_{s,t} \ge init_station_{s} \qquad \forall s \in \mathcal{S}, t \in \mathcal{T} \tag{init_station_status}
\end{align}